% !PRINTED 215
\weberpage{215}
%<<CONT
und die beiden Grundformen werden:
\weq{7}{\begin{aligned}
\Phi_1 &= x_1^4 + 2\sqrt{-3}\,x_1^2 x_2^2 + x_2^4 \\
\Phi_2 &= x_1^4 - 2\sqrt{-3}\,x_1^2 x_2^2 + x_2^4.
\end{aligned}}

Zwischen den drei Grundformen $f, \Phi_1, \Phi_2$ kann man eine
Relation herleiten, wenn man $x_1, x_2$ eliminirt. Man findet aus (7)
\[
2(x_1^4+x_2^4) = \Phi_1 + \Phi_2
\]
\[
4\sqrt{-3}\,x_1^2 x_2^2 = \Phi_1 - \Phi_2,
\]
und aus (1)
\[
f^2 = x_1^2 x_2^2 (x_1^4+x_2^4)^2 - 4x_1^6 x_2^6,
\]
woraus man leicht berechnet
\weq{8}{12\sqrt{-3}\,f^2 = \Phi_1^3 - \Phi_2^3.}

Die Substitutionen $\Theta, \psi, \chi$ erhalten, wenn sie mit der Determinante
 1 dargestellt werden, den Ausdruck:
\weq{9}{\begin{pmatrix} i, & 0 \\ 0, & -i \end{pmatrix},\quad
\begin{pmatrix} 0, & i \\ i, & 0 \end{pmatrix},\quad
\begin{pmatrix} \dfrac{1-i}{2}, & \dfrac{1-i}{2} \\[6pt] -\dfrac{1+i}{2}, & \dfrac{1+i}{2} \end{pmatrix},}
woraus nach den im \S.\,50 aufgestellten Bedingungen folgt, dass
die entsprechende Gruppe orthogonaler ternärer Substitutionen
reell ist.

Wendet man die Substitutionen (9) mit der Determinante 1
auf die Grundform $f(x_1, x_2)$ an, so ergiebt eine sehr einfache
\gesperrt{Rechnung, dass die Grundform $f$ eine absolute Invariante}
der binären Gruppe $G$ ist.

Dieselbe Eigenschaft hat auch die \gesperrt{Hesse'sche} Determinante
von $f$:
\[
\begin{aligned}
H &= -\frac{1}{25}\left[f''(x_1,x_1)\,f''(x_2,x_2) - f''(x_1,x_2)^2\right] \\
&= (x_1^4+x_2^4)^2 + 12\,x_1^4 x_2^4,
\end{aligned}
\]
die nichts Anderes ist, als das Product der beiden Functionen
$\Phi_1, \Phi_2$; während die Functionen $\Phi_1$ und $\Phi_2$ selbst bei der dritten
der Substitutionen (9) eine dritte Einheitswurzel als Factor
annehmen.
